\documentclass[10pt,a4paper]{amsart}
\usepackage[margin=19mm]{geometry}
\usepackage{amsmath,amssymb,mathtools}
\usepackage[scr=rsfs,cal=euler]{mathalfa}
\usepackage{xfrac}
\usepackage[unicode,hidelinks,bookmarks=false]{hyperref}
\newcommand{\ie}{i.e.\ }
\newcommand{\cf}{cf.\ }
\newcommand{\resp}{respectively\ }
\newcommand{\Subsection}{\subsection}
\providecommand{\nobreakdash}{\nobreak\hbox{-}}
\setlength{\emergencystretch}{2em}
\multlinegap0pt
\usepackage{amssymb}
\usepackage{enumerate}
\usepackage{xcolor}
\usepackage{tikz}
\usepackage{stackengine}
\newtheorem{theorem}{Theorem}
\newcommand{\B}{\mathcal{B}}
\title{\small A generalization of the inverse mapping theorem \\ \small in infinite dimensions}
\author{Sajjad Lakzian}
\date{Source: arXiv:2604.10002v1 (CC BY 4.0)}
\begin{document}
\maketitle

\noindent\emph{Source and licence.} \small A generalization of the inverse mapping theorem \\ \small in infinite dimensions. Version arXiv:2604.10002v1 (CC BY 4.0), distributed by arXiv under the Creative Commons Attribution 4.0 International licence, http://creativecommons.org/licenses/by/4.0/, as recorded in the arXiv licence metadata for this exact version and shown on the abstract page https://arxiv.org/abs/2604.10002v1. Full licence notice: \texttt{licenses/arxiv-2604.10002-CC-BY-4.0.txt}.\par\smallskip

\noindent\textit{Editorial scope.} Counted unit: the article's theorem theorem (source0.tex lines 958--963). The article's complete original proof of it is delivered with it (source0.tex lines 965--982, one \texttt{proof} environment of 18 lines). No mathematics is written, repaired or re-derived: the statement and the proof are the article's own lines, in the article's own notation, and no retained line is substituted or re-flowed.  Retained uncounted as the source-local context the proof needs: lines 639--647; lines 712--714; lines 951--955.  Nothing was omitted from any retained span, so the package carries no omission record.  No retained line contains a citation, so the wrapper carries no bibliography and no bibliography entry.  No retained line contains a figure environment, an image-inclusion command or drawing code, so no image is delivered and the package declares no asset dependency.  Editorial formatting only, in addition: an AMS \texttt{amsart} preamble that restates the article's own declarations and macros used by the retained lines, namely the environments \texttt{theorem} on separate counters, together with the standard packages those lines need.  The retained environments print in this document as Theorem 1 in order of appearance, whereas the article prints the same items with the numbering of its own full document.  The complete native source of the article as distributed in the arXiv e-print for this exact version is bundled unchanged as \texttt{sources/198240/source0.tex}.
\begin{theorem}[Generalized IMT]\label{thrm:main-1}
Let $f: \mathsf{Dom}(f) \subset \B_1 \to \B_2$ be an everywhere differentiable map with invertible derivative. Then $f$ is a local diffeomorphism provided any of the following items hold
	\begin{enumerate}
		\item if $f$ satisfies strong ${\sf A}$ everywhere on all scales;
		\item If $\B_1$ is locally weakly compact and $f$ is weak-weak continuous and satisfies ${\sf A}$ everywhere on all scales;
		\item If $\B_1$ is locally weakly compact with ${\sf FPP}$ and $f$ satisfies ${\sf A}$ everywhere on all scales; 
		\item $\B_1$ is strictly convex with ${\sf FPP}$ and $f$ satisfies weak ${\sf A}$ everywhere on all scales. 
	\end{enumerate} 
\end{theorem}

\begin{theorem}[Generalized implicit function theorem]\label{thrm:main-2}
	Suppose $\B_1$ and $\B_2$ are two Banach spaces and $g:W \subset \B_1 \times \B_2 \to \B_2$ is a function which we will write it as $g(t,x)$. Suppose $g$ is everywhere differentiable and with invertible derivative $D_x g(t,x)$ with respect to the variable $x$. Assume $g$ and the domain $\B_1\times \B_2$ satisfy either of the hypotheses of Theorem~\ref{thrm:main-1}. Then, every non-empty level set $g^{-1}(c)$ can be locally uniquely written, around $(a,b) \in g^{-1}(c)$ as a graph of a differentiable function $h: U_a \subset \B_1 \to \B_2$. 
\end{theorem}

\begin{align}\label{eq:abs-ode}
\begin{cases} u'(t) = f(t,u(t)), \quad f \in \mathcal{H}\\
	u(0) = b\in \B_1
\end{cases}.
\end{align}

\begin{theorem}
	The abstract pde system \eqref{eq:abs-ode} has a unique solution locally given by
	\[
	g(t,u(t)) = c = g(0,b).
	\]
\end{theorem}

\begin{proof}[\footnotesize \textbf{Proof}]
\par Let us rewrite the PDE \ref{eq:abs-ode}. From
\[
u'(t) = \left[D_xg(t,u(t))\right]^{-1} D_tg(t,u(t)),
\]
(this is a pde since $t$ comes from a multi-dimensional space)
we deduce
\[
D_xg(t,u(t))\circ u'(t) - D_tg(t,u(t)) = 0.
\]
which by virtue of the chain rule, the latter is just
\[
D_t\left( g(t,u(t))  \right) = 0.
\]
Thus any possible solution lies on a the level set $g^{-1}(c)$. 

\par By the IFT proven in Theorem~\ref{thrm:main-2}, level set $g^{-1}(c)$ around any of its points is locally the graph of a unique function $t \mapsto u(t)$. It is then obvious that $u(t)$ solves the system \eqref{eq:abs-ode} and is unique since by IFT, locally it is uniquely determined. 
\end{proof}

\end{document}
